\subsubsection{Introduction to resipy}

\paragraph{Overview}

\textbf{resipy} (also called \textit{ResIPy}) is a Python wrapper around the robust R2 / cR2 / R3t inversion codes for 2-D and 3-D electrical resistivity tomography (ERT).

\begin{framed}
\textbf{Learning Objectives}\\
\begin{itemize}
\item Understand the ERT data-acquisition workflow
\item Create and configure a resipy \texttt{Project}
\item Import field data and visualise a pseudosection
\item Run a basic regularised inversion
\item Interpret an inverted resistivity section
\end{itemize}
\end{framed}


\bigskip
\centerline{\rule{13cm}{0.4pt}}
\bigskip

\paragraph{Installation Check}

\begin{verbatim}
import numpy as np
import matplotlib.pyplot as plt
from resipy import Project

print("resipy imported successfully!")
\end{verbatim}

\begin{framed}
\textbf{Wine warning on Linux/macOS}\\
You may see \texttt{Warning: Wine is not installed!} --- this is expected if Wine is absent.
Data loading and visualisation work without Wine; only the inversion binary (R2.exe) needs it.
Install with \texttt{sudo apt install wine-stable} on Debian/Ubuntu to enable full inversion.
\end{framed}


\bigskip
\centerline{\rule{13cm}{0.4pt}}
\bigskip

\paragraph{Creating a Project}

\begin{verbatim}
# Create a 2 -D resistivity project
k = Project(typ='R2')
print("Project created!")
\end{verbatim}

The \texttt{typ} argument selects the inversion engine:

\bigskip\noindent
\begin{tabular}{p{\dimexpr 0.333\linewidth-2\tabcolsep}p{\dimexpr 0.333\linewidth-2\tabcolsep}p{\dimexpr 0.333\linewidth-2\tabcolsep}}
\toprule
\texttt{typ} & Engine & Use case \\
\hline
\texttt{'R2'} & R2.exe & 2-D resistivity \\
\texttt{'cR2'} & cR2.exe & 2-D complex resistivity (IP) \\
\texttt{'R3t'} & R3t.exe & 3-D resistivity \\
\texttt{'cR3t'} & cR3t.exe & 3-D complex resistivity \\
\bottomrule
\end{tabular}

\bigskip
\bigskip
\centerline{\rule{13cm}{0.4pt}}
\bigskip

\paragraph{Importing Survey Data}

\begin{verbatim}
# Load a survey data file (ABMN format)
# k.createSurvey('data/sample_resistivity/survey.dat')

# For this demo, we inspect the expected data format
print("Expected column format (.dat file):")
print("  A    B    M    N    Resistance( \Omega)")
print("  0    9    3    6    12.45")
print("  0   12    4    8     9.21")
print("  ...")
\end{verbatim}

\begin{framed}
\textbf{Data formats}\\
resipy supports many common ERT data formats including:
\textbf{Syscal} (IRIS), \textbf{Res2Dinv}, \textbf{ABEM LS}, \textbf{Campus Tigre}, \textbf{Protocol} and more.
Pass \texttt{ftype='Syscal'} (or other) to \texttt{createSurvey()} for automatic parsing.
\end{framed}


\bigskip
\centerline{\rule{13cm}{0.4pt}}
\bigskip

\paragraph{Pseudosection Visualisation}

A \textbf{pseudosection} plots apparent resistivity against electrode mid-point and pseudo-depth.
It is a qualitative picture of the data --- not a true depth section.

\begin{verbatim}
# After loading data:
# k.showPseudoSection()

# Simulated pseudosection for illustration
np.random.seed(42)
n_levels = 8
n_points = 40

fig, ax = plt.subplots(figsize=(11, 5))

x_mid  = np.linspace(0, 47, n_points)
depths = np.arange(1, n_levels + 1) * 1.5

cmap = plt.cm.RdYlBu_r
norm = plt.Normalize(vmin=10, vmax=500)

for depth in depths:
    rho_a = np.exp(np.random.normal(loc=4.5, scale=0.4, size=n_points))
    sc = ax.scatter(x_mid, -np.full_like(x_mid, depth),
                    c=rho_a, cmap=cmap, norm=norm, s=60, marker='v')

cb = fig.colorbar(sc, ax=ax, pad=0.02)
cb.set_label("Apparent Resistivity ( \Omega·m)")
ax.set_xlabel("Mid -point position (m)")
ax.set_ylabel("Pseudo -depth (m)")
ax.set_title("Synthetic Pseudosection — Wenner Array")
ax.grid(True, alpha=0.2)
plt.tight_layout()
plt.show()
\end{verbatim}


\bigskip
\centerline{\rule{13cm}{0.4pt}}
\bigskip

\paragraph{Running an Inversion}

\begin{verbatim}
# After creating the survey and mesh:
#
# k.createMesh(typ='trian')        # triangular mesh
# k.param['a_wgt'] = 0.01         # regularisation weight
# k.invert()                       # run R2
# k.showResults()                  # display inverted section
#
# The inversion minimises:
#   \Phi = ‖Wd(d_obs - d_pred)‖² + \lambda‖Wm m‖²
# where \lambda is the regularisation parameter (trade -off).

print("Inversion parameters overview:")
params = {
    "a_wgt"    : "Absolute data weight (noise floor)",
    "b_wgt"    : "Relative data weight",
    "alpha_s"  : "Smoothness regularisation weight",
    "max_iter" : "Maximum inversion iterations",
}
for k_p, desc in params.items():
    print(f"  {k_p:12s}  \rightarrow  {desc}")
\end{verbatim}


\bigskip
\centerline{\rule{13cm}{0.4pt}}
\bigskip

\paragraph{Synthetic Inverted Section}

\begin{verbatim}
from matplotlib.colors import LogNorm

# Synthetic 2 -D resistivity model for illustration
nx, nz = 80, 30
x = np.linspace(0, 47, nx)
z = np.linspace(0, -15, nz)
X, Z = np.meshgrid(x, z)

# Background + conductive anomaly + resistive layer
rho = np.full((nz, nx), 150.0)
rho[Z > -8] = 80      # shallow moist layer
rho[(Z > -12) & (Z < -8) & (X > 15) & (X < 30)] = 12   # conductive body
rho[Z < -13] = 800    # resistive bedrock

fig, ax = plt.subplots(figsize=(11, 4))
im = ax.pcolormesh(X, Z, rho, cmap='RdYlBu_r', norm=LogNorm(vmin=10, vmax=1000), shading='auto')
cb = fig.colorbar(im, ax=ax, pad=0.02)
cb.set_label("Resistivity ( \Omega·m)")
ax.set_xlabel("Distance (m)")
ax.set_ylabel("Depth (m)")
ax.set_title("Synthetic Inverted Resistivity Section")
plt.tight_layout()
plt.show()
\end{verbatim}


\bigskip
\centerline{\rule{13cm}{0.4pt}}
\bigskip

\paragraph{Interpretation Guide}

\begin{table}
\centering
\begin{tabular}{p{\dimexpr 0.333\linewidth-2\tabcolsep}p{\dimexpr 0.333\linewidth-2\tabcolsep}p{\dimexpr 0.333\linewidth-2\tabcolsep}}
\toprule
Feature & Resistivity & Possible interpretation \\
\hline
Shallow low-$\rho$ zone & \textless  30 $\Omega$·m & Moist / clay-rich topsoil \\
Deep conductive body & 5--20 $\Omega$·m & Saline groundwater, leachate plume \\
Resistive layer & \begin{quote}
500 $\Omega$·m
\end{quote}

 & Dry sand, gravel, bedrock \\
Lateral contrast & --- & Lithological boundary or fault \\
\bottomrule
\end{tabular}
\end{table}


\bigskip
\centerline{\rule{13cm}{0.4pt}}
\bigskip

\begin{framed}
\textbf{Summary}\\
You have learned:

\begin{itemize}
\item How to create a \textbf{resipy} \texttt{Project} and load ERT data
\item What a pseudosection shows (and its limitations)
\item How regularised inversion works conceptually
\item Basic interpretation of an inverted resistivity section
\end{itemize}

Next: \href{01\_emagpy\_introduction.md}{emagpy Introduction} --- process \textbf{electromagnetic} data with emagpy.
\end{framed}